\documentclass[a4paper,11pt]{article}
\usepackage[a4paper,top=0.56in,bottom=0.55in,left=0.65in,right=0.65in]{geometry}
\usepackage[T1]{fontenc}
\usepackage[utf8]{inputenc}
\usepackage[sfdefault]{roboto}
\usepackage{microtype}
\usepackage{xcolor}
\usepackage{enumitem}
\usepackage{titlesec}
\usepackage[hidelinks]{hyperref}
\usepackage{tabularx}
\usepackage{ragged2e}
\input{glyphtounicode}
\pdfgentounicode=1
\definecolor{accent}{RGB}{35,61,89}
\pagestyle{empty}
\setlength{\parindent}{0pt}
\setlength{\parskip}{0pt}
\setlength{\emergencystretch}{1em}
\urlstyle{same}
\titleformat{\section}{\bfseries\normalsize\color{accent}}{}{0pt}{}[\vspace{-3pt}\color{accent}\titlerule]
\titlespacing*{\section}{0pt}{7.5pt}{3.5pt}
\newcommand{\position}[4]{\textbf{#1}\hfill\textbf{#2}\\[-1pt]\textit{#3}\hfill\textit{#4}\vspace{-3pt}}
\newcommand{\proj}[2]{\textbf{#1}\hfill{\footnotesize #2}\vspace{-4pt}}
\newenvironment{tightbullets}{\begin{itemize}[leftmargin=1.25em,label=\textbullet,topsep=2pt,itemsep=2pt,parsep=0pt,partopsep=0pt]}{\end{itemize}\vspace{-4pt}}
\newcommand{\skill}[2]{\textbf{#1:} #2\\[1pt]}
\begin{document}
\begin{center}
{\fontsize{17.3}{20}\selectfont\bfseries Devanshu Kumar Ranjan}\\[2pt]
{\small\bfseries AI ENGINEER | AI APPLICATION SECURITY | SECURITY AUTOMATION}\\[4pt]
{\footnotesize Bilaspur, Himachal Pradesh, India \quad | \quad +91-6206938457 \quad | \quad \href{mailto:devanshukranjan@gmail.com}{devanshukranjan@gmail.com}}\\[2pt]
{\footnotesize \href{https://linkedin.com/in/devanshu-kumar-ranjan-06399221b}{linkedin.com/in/devanshu-kumar-ranjan-06399221b} \quad | \quad \href{https://github.com/devanshukranjan}{github.com/devanshukranjan}}\\[2pt]
{\footnotesize \href{https://devanshukranjan.github.io/AI-Cybersecurity-Portfolio/}{devanshukranjan.github.io/AI-Cybersecurity-Portfolio}}
\end{center}
\vspace{-6pt}
\section{Professional Summary}
\small Production AI Engineer and NIT Patna M.Tech Cyber Security graduate focused on \textbf{evidence-grounded agentic AI}, prompt-injection safeguards and secure Python services. Built a LangGraph RCA investigator with 30 completed diagnostics, read-only PostgreSQL access and AI reasoning; developed Splunk and ML-IDS labs.\normalsize

\section{Professional Experience}
\position{KAS Global Commerce Inc.}{Bilaspur, HP}{AI Engineer}{Jul 2026 -- Present}
\begin{tightbullets}
\item Built a production \textbf{LangGraph RCA investigator} with Python evidence retrieval using read-only PostgreSQL table access and AI-generated causal explanations across \textbf{30 completed diagnostics} in advertising, orders, business and inventory for seller questions.
\item Integrated \textbf{prompt-injection safeguards}, validated input, statistical comparison checks and typed evidence-backed reporting to constrain AI conclusions and recommendations.
\item Lowered PostgreSQL analytics latency from \textbf{over 60 seconds to under 7 seconds (88\%+)} using parallel query execution and indexing.
\item Automated monthly catalog ingestion and multimodal insight generation, with recovery handling, across \textbf{9 global marketplaces}.
\end{tightbullets}
\vspace{2pt}
\position{KAS Global Commerce Inc.}{Bilaspur, HP}{AI Research Intern}{Sep 2025 -- Jun 2026}
\begin{tightbullets}
\item Built a two-stage Zyte/MongoDB catalog pipeline, reducing extraction from \textbf{2--3 months to 1.5 weeks}; deployed taxonomy vector search using \texttt{pgvector} with \textbf{sub-50 ms retrieval}.
\item Built \textbf{Gemini 2.5 Flash} text/image category intelligence from scraped Amazon/ONDC listings for competitor patterns, buyer preferences and listing guidance; stored month-versioned insights in marketplace-specific PostgreSQL \textbf{JSONB} columns.
\item Developed \textbf{3 FastAPI microservices}; \texttt{/api/v1/insights} served precomputed category intelligence to the MSC IQ seller assistant.
\end{tightbullets}

\section{Security and AI Research Projects}
\proj{\href{https://github.com/devanshukranjan/AI-Cybersecurity-Portfolio/tree/main/siem-investigations/splunk-siem-lab}{Splunk SIEM: Security Monitoring and Detection Lab}}{Apr -- May 2026}
\begin{tightbullets}
\item Built a \textbf{3-VM SOC lab}; collected Linux/Apache logs, simulated Hydra/Nikto events, and wrote SPL detection searches with automated threshold alerts.
\end{tightbullets}
\proj{\href{https://github.com/devanshukranjan/AI-Cybersecurity-Portfolio/tree/main/ai-intrusion-detection/hermes-ids}{Hermes: AI-Assisted Intrusion Detection Prototype}}{May -- Jun 2025}
\begin{tightbullets}
\item Combined Random Forest (\textbf{98.6\% accuracy}), Tab2Image CNN, and autoencoder anomaly detection on network traffic datasets; built a Tkinter security analyst console.
\end{tightbullets}
\proj{Privacy-Preserving Federated Sequential Recommendation (M.Tech thesis)}{Aug 2025 -- May 2026}
\begin{tightbullets}
\item Developed \textbf{LDP-SeqMask} using client-side Laplace perturbation and view-based contrastive learning to protect interaction sequences in federated recommendation.
\end{tightbullets}

\section{Technical Skills}
{\small
\skill{AI Application Engineering}{Python, FastAPI, LangGraph, RAG, AI agents, LLM prompting, Gemini 2.5 Flash, PostgreSQL, pgvector}
\skill{AI Security and Privacy}{Prompt-injection input safeguards (production), local differential privacy, federated learning, secure API design fundamentals}
\skill{Detection and Investigation (Labs)}{Splunk, SPL, auth/Apache logs, Wireshark, network traffic analysis, IDS/ML, Linux, Git}
}
\section{Education and Certification}
{\small\textbf{NIT Patna} --- M.Tech, Cyber Security (CSE), \textbf{9.06/10} \hfill 2024 -- 2026\\[1pt]
\textbf{Graphic Era University} --- B.Tech, CSE, \textbf{8.83/10} \hfill 2019 -- 2023\\[1pt]
\textbf{Certification:} Google Cybersecurity Specialization (Coursera), 2025}
\end{document}
